\medskip

Deux avions sont en approche d'un aéroport.

On munit l'espace d'un repère orthonormé \Oijk dont J'origine O est le pied de la
tour de contrôle, et le sol est le plan $P_0$ d'équation $z = 0$.

L'unité des axes correspond à 1 km.

On modélise les avions par des points.

\begin{center}
\psset{unit=1cm,arrowsize=2pt 3}
\begin{pspicture}(-6,-1.8)(4.7,5.5)

\psline{->}(-6,0.3)(4.7,-0.1)
\psline{->}(0,-1.8)(0,5.5)
\psline{->}(-3,-1)(3,1)
\multido{\n=-1.3+1.3,\nb=-2+2}{5}{\rput(0,\n){\small $\bullet$}\uput[l](0,\n){\small \nb}}
\multido{\n=-4.8+1.2,\na=0.25 +-0.05,\nb=-8+2}{8}{\rput(\n,\na){\small $\bullet$}\uput[d](\n,\na){\small \nb}}
\multido{\n=-2.4+0.6,\na=-0.8 +0.2,\nb=-8+2}{9}{\rput(\n,\na){\small $\bullet$}\uput[ul](\n,\na){\small \nb}}
\psline[linewidth=1.5pt](-4.1,5.5)(-1.7,-0.3)\psline[linewidth=1.5pt,linestyle=dashed](-1.7,-0.3)(-1.4,-1)
\psline[linewidth=1.5pt](-6,5.5)(1,-0.4)\psline[linewidth=1.5pt,linestyle=dashed](1,-0.4)(1.8,-1)
\uput[d](4.6,-0.15){$x$}\uput[d](0,5.3){$z$}\uput[ul](3,1){$y$}
\uput[ur](-3.5,4){$d_{\text{A}}$}\uput[dl](-4.2,4){$d_{\text{B}}$}
\rput(-3.9,5){$\bullet$}\rput(-4.8,4.5){$\bullet$}
\uput[ur](-3.9,4.95){A}\uput[dl](-4.8,4.5){B}
\end{pspicture}
\end{center}

L'avion Alpha transmet à la tour sa position en A$(-7~;~ 1~;~7)$ et sa trajectoire est dirigée par le vecteur $\vect{u}\begin{pmatrix}2\\-1\\-3\end{pmatrix}$.

L'avion Bêta transmet une trajectoire définie par la droite $d_{\text{B}}$ passant par le point B dont une représentation paramétrique est :

\[\left\{\begin{array}{l c l}
x &=&-11 + 5t\\
y &=&-5 + \phantom{5}t\\
z &=&11 - 4t \end{array}\right.\:\text{où}\: t\:\text{décrit}\:\: \R\]

\medskip

\begin{enumerate}
\item S'il ne dévie pas de sa trajectoire, déterminer les coordonnées du point S en lequel l'avion Bêta touchera le sol.
\item 
	\begin{enumerate}
		\item Déterminer une représentation paramétrique de la droite $d_{\text{A}}$ caractérisant la trajectoire de l'avion Alpha.
		\item Les deux avions peuvent-ils entrer en collision ?
	\end{enumerate}
\item 
	\begin{enumerate}
		\item Démontrer que l'avion Alpha passe par la position E$(-3~;~-1~;~1)$.
		\item Justifier qu'une équation cartésienne du plan $P_{\text{E}}$ passant par E et perpendiculaire à la droite $d_{\text{A}}$ est:

\[2x - y - 3z + 8 = 0.\]
		\item Vérifier que le point F$(-1~;~-3~;~3)$ est le point d'intersection du plan $P_{\text{E}}$ et de la droite $d_{\text{B}}$.
		\item Calculer la valeur exacte de la distance EF, puis vérifier que cela correspond à une distance de \np{3464} m, à 1 m près.
	\end{enumerate}
\item La réglementation aérienne stipule que deux avions en approche doivent être à tout instant à au moins 3 milles nautiques l'un de l'autre (1 mille nautique vaut \np{1852}~m).

Si les avions Alpha et Bêta sont respectivement en E et F au même instant, leur distance de sécurité est-elle respectée ?
\end{enumerate}

\bigskip
